% !PRINTED 432
\weberpage{432}
%<<CONT
\[
\begin{array}{c|cccccccccccc}
I & 0 & 1 & 2 & 3 & 4 & 5 & 6 & 7 & 8 & 9 & 10 & 11 \\ \hline
N & 1 & 2 & 4 & 8 & 3 & 6 & 12 & 11 & 9 & 5 & 10 & 7
\end{array}
\]

\[
\begin{array}{c|cccccccccccc}
N & 1 & 2 & 3 & 4 & 5 & 6 & 7 & 8 & 9 & 10 & 11 & 12 \\ \hline
I & 0 & 1 & 4 & 2 & 9 & 5 & 11 & 3 & 8 & 10 & 7 & 6
\end{array}.
\]

In the Canon Arithmeticus of \gesperrt{Jacobi} (1839) index tables
are compiled for all primes in the first thousand. On a smaller
scale, after computations of \gesperrt{Ostrogradsky}, in the ``Theory of
Congruences'' of \gesperrt{Tschebyscheff} (German by
\gesperrt{Schapira}, Berlin 1889).

It is still to be remarked that the index of 1 is always equal to 0,
and that the index of $-1$ (or of $n-1$) is always equal to
$\dfrac{n-1}{2}$. For since $(-1)^2$ has the index 0 or $n-1$, $-1$
must have the index $\dfrac{n-1}{2}$, since it does not have the
index 0, and twice its index is equal to 0 or $n-1$; in formulae:
\weq{21}{\operatorname{ind} 1 \equiv 0, \quad \operatorname{ind}(-1) \equiv \frac{n-1}{2} \pmod{n-1}.}

\begin{center}
\S.\,137.

\gesperrt{Multiplication and division of the trigonometric functions.}
\end{center}

We have already pointed out the connection that exists between the
theory of the roots of unity and the geometric problem of the
division of the circumference of the circle. The generalisation of
this problem leads to the division of an arbitrary angle. This
problem too leads to algebraic equations, of which we have already
used the one for trisection earlier, for the general solution of the
cubic equations (\S.\,112).

We now begin with the setting up of the general formulae.

By the theorem of \gesperrt{Moivre} (\S.\,33), if $n$ is an
arbitrary positive integer and $\varphi$ an arbitrary angle, we have
%CONT>>
